% !TEX root = ../manuscript.tex

\section{Discussion}

\todo{Open with three quantitative findings: the dominant geologic controls,
the geologic settings in which stochastic realizations change the containment
class, and the observed performance of the representative-fault benchmark.}

\subsection{Why end-to-end uncertainty propagation is needed}

The contribution is not simply the connection of several numerical tools.
The framework preserves the chain from fault--seal geology to plausible local
fault-core architectures, physically consistent
$(\mathbf{K},\phi,P_c,k_r)$ properties, full-fault fields, and
reservoir-scale outcome distributions. This common chain permits direct
comparison of broad geologic effects, residual within-geology stochastic
effects, and representative versus ensemble predictions without replacing
joint fault properties by independently averaged components.

\subsection{When is stochastic fault modeling needed?}

The required level of fault representation should be tied to the stability of
the storage conclusion. Where stochastic realizations are consistently
contained and the medoid agrees with the ensemble, a representative fault may
be adequate for initial assessment. Mixed containment outcomes indicate that
an ensemble is needed and that additional fault characterization or monitoring
has greater value. Where all realizations migrate substantially, further
sampling may refine the magnitude without changing the broad site or
injection-management conclusion. \todo{Replace these conditional statements
with a compact result-supported guide after the enriched simulations are
analyzed.}

\subsection{Relation to previous uncertainty-propagation studies}

Pettersson et al.\ demonstrate stochastic upscaling, preserve dependence among
multiphase fault properties, and propagate those properties to field-scale
leakage distributions using a one-dimensional fine-scale fault description
\cite{Pettersson2025}. Our framework extends the scientific question to a
broad fault--seal design space and a spatially variable three-dimensional
fault: where does additional stochastic detail change the storage conclusion,
and where does it not? This distinction is more useful than assuming either
that one representative fault is universally sufficient or that every site
requires the same ensemble size.

These findings should be interpreted within the scope of the modeled design.
The 162 scenarios span a balanced geologic design space but do not carry
field-calibrated occurrence probabilities. PREDICT realizations are
model-based samples conditional on the prescribed geology, not observations
of the real fault. Full-fault stochastic fields assume conditional
independence after conditioning on geology and throw window because available
data do not constrain a coarse-scale spatial covariance model; they therefore
omit uncertainty in unresolved large-scale fault continuity. The enriched
ensembles support common and moderately uncommon outcomes, not precise
rare-event probabilities. Results also remain conditional on the reservoir
model, constitutive assumptions, injection scenario, and exclusion of
geomechanical reactivation.
